\scanpage{242}
\mspage{188}
% archive sheet number: 241
même un $W$-fibré, ce qui signifierait que
\[
{}_{s'\backslash}X\longrightarrow{}_{s\backslash}X
\qquad\text{associé à}\qquad s\mathrel{\unsure{\text{---}|\longrightarrow}}s'\quad\text{dans }S
\]
serait toujours dans $W$. Posant $Z'={}_{s'\backslash}S$, $Z={}_{s\backslash}S$,
cela signifierait que si on a un \inserted{co-}\uline{\unsure{isom|isom.} local}
\[
 i:Z'\longrightarrow Z
\]
où $Z',Z$ ont \deleted{des} \inserted{$\unsure{z'_0,z_0|z'_0\;z_0}$} éléments initiaux, et si
$X$ est cofibrée sur $Z$, $X'$ son image inverse sur
$Z'$, alors $X'\mathrel{\unsure{\text{---}|\longrightarrow}}X$ est dans $W$. \deleted{\deleted{\unsure{Voici un}
contre-exemple, soit} $Z=\Delta^2=\{0,1,2\}$, $Z'=\Delta^1=\{0,1\}$,
$i$ l'inclusion évidente (identique), qui est une
immersion fermée donc un coisom. local. La
donnée de $X$ cofibrée et \unsure{coscindée} sur $\deleted{S}\inserted{Z}=\Delta^2$
équivaut à une suite de deux foncteurs
\[
X_0\mathrel{\unsure{\text{---}|\longrightarrow}}X_1\mathrel{\unsure{\text{---}|\longrightarrow}}X_2
\]
(où $X_i=$ fibre de $X$ en $i$), $X'$ \unsure{estant|étant} \unsure{décrit|décrite}
alors par $X_0\longrightarrow X_1$. Comme $X$ $W$-propre \uline{,,} sur
$S$, a un élément final $e=2$, l'inclusion $X_e\mathrel{\unsure{\text{---}|\longrightarrow}}X_{/e}=X$
est $\in W$, donc dans $\mathrm{Hot}_W$ \deleted{$X$ s'identifie à} $X_e=X_2$
\deleted{de $=$ $X'$ s'identifie à $X_1$}} \unsure{toute fois|toutefois}, si $Z$ a \inserted{aussi} un
élément \uline{final} $e$, (comme c'était le cas de départ
pour $\Delta^1$) alors $\deleted{\unsure{X\simeq X_e}}$ $Z'\simeq{}_{z_0\backslash}Z$ en a aussi
un $e'$, qui s'envoie sur $e$, et \deleted{\unsure{de même}} on aura
le diagramme commutatif
\diagram{\begin{tikzcd}[ampersand replacement=\&,column sep=small,row sep=small]
X'\arrow[r,"j"]\&X\\
X'_{e'}\arrow[u,"W"]\arrow[r,"\sim"]\&X_e\arrow[u,"W"']
\end{tikzcd}}
, d'où $j\in W$. Donc pour faire un contreexemple,
